\section{Q\&A：成本、目标函数与价值多元}

准备好的 slides 在 52:40 左右结束，后续问答反复回到早先图表，却加入了课堂最重要的限制条件。本章只保留新增 spoken evidence，不重复插图。

\subsection{Descriptive Morality 与 Prescription 再次分开}

学生追问 Delphi 是否把现有判断当成应该遵守的道德。讲者回答，研究目标首先是建模 descriptive judgments，并承认其中包含 bias；normative endorsement 需要额外原则和社会程序。她还把 ChatGPT 描述为更强的 generation model，而不是已经完成的 knowledge model。

\teachervoice{Q\&A 中的关键澄清是：能预测“人们会怎么说”不等于知道“应该怎么做”。即使 prediction accuracy 很高，系统仍需展示数据群体、分歧分布与拒答条件。}

\subsection{Culture、Dataset 与 Distribution of Opinions}

道德数据来自不同年代、国家、平台和 annotation instruction。讲者指出 Delphi 有 Western bias，也提到不同 moral datasets 的规模与覆盖不等。更合理的目标不是在所有争议问题上输出单一答案，而是区分 unacceptable harm 与 legitimate disagreement，并表示 opinion distribution。

\begin{knowledgebox}{Pluralism 的机器学习表示}
单标签 classification 会把分歧压成多数票。更完整的输出可以包含 label distribution、群体条件、confidence、rationale clusters 与 contested flag。这样仍不能解决规范冲突，但至少不会把 disagreement 隐藏成确定性。
\end{knowledgebox}

\subsection{Complexity：Maieutic 与 Hybrid 都很慢}

递归 explanation expansion 会多次调用 LM，weighted MaxSAT 需要在图上求解，hybrid 还加入 retrieval 和 principle checking。讲者直接承认 Maieutic prompting 与 Delphi hybrid 都慢。系统验收因此必须报告 token calls、tree width/depth、solver time 与 end-to-end latency，而不能只报 accuracy。

\subsection{COMET 与 GPT-3 的比较不是 Architecture Ablation}

问答再次确认 COMET 使用专门 ATOMIC 监督，GPT-3 使用 few-shot prompt，critical teacher 又经过 critic filtering。结果可以支持“任务专门化与数据质量有价值”，却不能单独估计 parameter count、pretraining corpus 或 architecture 的因果贡献。

\subsection{为什么人类也需要 Critic}

讲者用学习过程解释 critic：人类并不是第一次想到什么就永远相信什么，而会通过反馈、反思和社会纠正形成 internal critic。这个类比提供设计直觉，但机器 critic 的训练分布仍需显式审计，不能因为“像人”就获得可信度。

\subsection{Refusal 不能取代 Context Understanding}

有学生问系统是否应完全回避 moral / ethical questions。讲者认为直接 moral advice 可以限制，但系统仍必须识别骚扰、危险、仇恨或求助语境；否则只靠拒答会让安全组件看不见真正的 harmful context。工程上需要 answer policy 与 understanding model 分离。

\begin{importantbox}{Abstention 也是一项需要训练和评估的能力}
好的拒答不是遇到关键词就停止，而是能说明不确定性、识别高风险、提供安全替代、记录触发原因，并把需要权威判断的问题升级给人类。abstention 应报告 false refusal 与 unsafe answer 两类错误。
\end{importantbox}

\subsection{Next-Token、Preference 与 Knowledge 是不同目标}

问答末尾回到本讲最深的主张：next-token prediction 学语言延续，RLHF / preference tuning 学人类偏好，knowledge model 学世界关系与可查询推断。三者会共享表示并互相促进，却不是同一个 objective。commonsense 可能从 LM 中 emerge，但 emerge 不等于稳定、可校准或可审计。

\teachervoice{讲者说 adversarial commonsense 的空间几乎无限，人类面对第一次听见的荒谬组合也能依靠 conceptual understanding 判断，而 Transformer 更偏向预测接下来哪个词。这是她坚持“language model 不等于 knowledge model”的最终理由。}

\subsection{Humanities 不是 External Reviewer，而是共同设计者}

最后的 value pluralism 问题把责任扩展到 humanities at large。哪些差异属于可接受分歧、哪些伤害应被拒绝，不能只由 AI researchers 决定。讲者用法律作为类比：规则是社会协商形成的人类 artifact，模型可以辅助表示和检查，但合法性来自制度参与。

\subsection{本章小结}

Q\&A 把 prepared talk 的技术方案重新放回资源、目标和治理。Maieutic / hybrid 有明显计算成本；specialist 与 generalist comparison 有混杂因素；拒答不等于理解；价值多元需要分布式表示与跨学科制度参与。

